% Generated from paper/full_draft. Edit the Markdown source or migration script.

\section{Experimental Setup}
\label{sec:05-experimental-setup:experimental-setup}

\subsection{Datasets}
\label{sec:05-experimental-setup:datasets}

The experiments use several datasets, but they have different roles. The paper does not treat all datasets as equal evidence for the main claim.

\begin{table*}[t]
\centering
\small
\caption{Dataset roles and evaluation guardrails}
\label{tab:1}
\begin{tabularx}{\textwidth}{XXXX}
\toprule
Dataset & Role & Paper use & Caveat \\
\midrule
LoTTE technology/search & Main vertical-domain retrieval benchmark & Main scale-up, token-quality frontier, geometry validation & No true corpus topic labels in processed qrels \\
PubMedQA & Feedback/manifold proof-of-concept & Shows trust feedback and local propagation can improve policy & GT is abstract-level context, not strict answer sentence \\
Banking77 & Intent/domain routing proxy & Shows strong feedback self-evolution and intent structure & Should not be mixed with evidence retrieval main table \\
eManual & Failure/limitation case & Shows duplicate text and strict chunk-id issues & Low strict recall does not prove geometry is absent \\
CUAD & Sparse smoke/stress case & Shows sparse legal-domain limitation & GT-anchored sample only, not full-corpus main evidence \\
\bottomrule
\end{tabularx}
\end{table*}

LoTTE technology/search is the main large-scale evidence benchmark. We evaluate nested corpus scales from 100k to 638k chunks with 596 test queries. CUAD and eManual are reported as limitation cases rather than main positive evidence.

\subsection{Baselines and Variants}
\label{sec:05-experimental-setup:baselines-and-variants}

The baseline family includes:

\begin{itemize}
\item BM25-only lexical retrieval.
\item Dense-only retrieval with \texttt{sentence-transformers/all-MiniLM-L6-v2}.
\item BM25 + dense hybrid retrieval using reciprocal-rank fusion.
\item Full multi-route IntentWeight.
\item Gated cost-aware IntentWeight.
\item Confidence-based final context IntentWeight.
\item Static geometry controls such as nearest-cluster routing.
\item Naive controls such as random or epsilon-greedy arm selection.
\end{itemize}

Dense-only retrieval is the primary quality baseline. The paper should avoid weak baseline framing: dense is strong and remains a required recall floor in the proposed method.

\subsection{Metrics}
\label{sec:05-experimental-setup:metrics}

Retrieval quality is measured with:

\begin{itemize}
\item $\mathrm{Hit@K}$: whether any ground-truth chunk appears in the top $K$.
\item $\mathrm{EvidenceRecall@K}$: fraction of all ground-truth chunks retrieved.
\item $\mathrm{MRR@K}$: reciprocal rank of the first relevant chunk.
\item $\mathrm{nDCG@K}$: binary relevance ranking quality.
\end{itemize}

For query $q_t$, let $G_t$ be the set of ground-truth chunks and let $R_t^K$ be the top-$K$ retrieved chunks. The query-level hit metric is:

\[
\mathrm{Hit@K}(q_t) =
\mathbb{1}\left[ R_t^K \cap G_t \neq \varnothing \right].
\]

The evidence-recall metric is:

\[
\mathrm{EvidenceRecall@K}(q_t) =
\frac{\left| R_t^K \cap G_t \right|}{\left|G_t\right|}.
\]

For mean reciprocal rank, let $\rho_t$ be the rank of the first relevant chunk within the top-$K$ list, or $\infty$ if no relevant chunk is retrieved:

\[
\mathrm{MRR@K}(q_t) =
\begin{cases}
\frac{1}{\rho_t}, & \rho_t \le K, \\
0, & \rho_t = \infty.
\end{cases}
\]

For binary relevance, let $\mathrm{rel}_{t,j} \in \{0,1\}$ denote whether the chunk at rank $j$ for query $q_t$ is relevant. We compute:

\[
\mathrm{DCG@K}(q_t) =
\sum_{j=1}^{K} \frac{\mathrm{rel}_{t,j}}{\log_2(j+1)},
\]

\[
\mathrm{nDCG@K}(q_t) =
\frac{\mathrm{DCG@K}(q_t)}{\mathrm{IDCG@K}(q_t)}.
\]

If $\mathrm{IDCG@K}(q_t)=0$, we set $\mathrm{nDCG@K}(q_t)=0$.

The main headline uses query-level $\mathrm{Hit@10}$. This choice reflects the target use case: retrieving at least one usable evidence chunk for RAG generation under a smaller context budget. It does not imply complete evidence collection. $\mathrm{EvidenceRecall@10}$ is reported separately where multi-evidence coverage is important.

Cost and efficiency are separated into three layers:

\begin{itemize}
\item Source candidate cost: number of candidates considered before final fusion.
\item Dense invocation rate: fraction of queries using the global dense path.
\item Final context tokens: token count of retrieved chunks sent to the generator.
\end{itemize}

The main cost result uses final context tokens. Source candidate cost and dense invocation rate are retrieval-stage diagnostics.

Let $C_t$ be the final retrieved context selected for query $q_t$ and let $\mathrm{tok}(d)$ be the token count of chunk $d$. The final context token cost is:

\[
\mathrm{Tokens}(C_t) =
\sum_{d \in C_t} \mathrm{tok}(d).
\]

\subsection{Geometry Diagnostics}
\label{sec:05-experimental-setup:geometry-diagnostics}

The geometry diagnostics are used to test whether the corpus has usable local structure for route control. They are diagnostic support for the piecewise relevance-manifold assumption, not mathematical proof of a manifold.

Let $E \in \mathbb{R}^{n \times p}$ be a sampled matrix of corpus chunk embeddings after mean-centering. Let $\lambda_1 \ge \lambda_2 \ge \cdots \ge \lambda_p \ge 0$ be the eigenvalues of the empirical covariance matrix. The explained variance retained by the first $m$ principal components is:

\[
\mathrm{PCAvar@m} =
\frac{\sum_{i=1}^{m} \lambda_i}
     {\sum_{i=1}^{p} \lambda_i}.
\]

The dimension needed to explain 90\% of the variance is:

\[
\mathrm{PCAdim90} =
\min \left\{m:
\frac{\sum_{i=1}^{m} \lambda_i}
     {\sum_{i=1}^{p} \lambda_i}
\ge 0.9
\right\}.
\]

For cluster diagnostics, let $y_t$ be the PCA/context representation of query $q_t$, let $\mu_c$ be the centroid of cluster $c$, and let $\ell_i$ be the cluster label of chunk $d_i$. The top-$K$ nearest clusters for query $q_t$ are:

\[
\mathcal{N}_K(q_t) =
\operatorname{TopK}_{c}
\left(\mu_c^\top y_t\right).
\]

Let $\mathcal{C}_t = \{\ell_i : d_i \in G_t\}$ be the set of clusters that contain ground-truth evidence for query $q_t$. The nearest-cluster hit metric is:

\[
\mathrm{NearestClusterHit@K} =
\frac{1}{|\mathcal{Q}_{GT}|}
\sum_{q_t \in \mathcal{Q}_{GT}}
\mathbb{1}
\left[
\mathcal{N}_K(q_t) \cap \mathcal{C}_t \neq \varnothing
\right],
\]

where $\mathcal{Q}_{GT}$ is the set of queries with at least one ground-truth chunk in the evaluated corpus.

Finally, let $R_{t,\mathrm{ctx}}^K$ be the top-$K$ chunks retrieved by inner product in the PCA/context space, and let $R_{t,\mathrm{dense}}^K$ be the top-$K$ chunks retrieved by dense embedding similarity. We define:

\[
\mathrm{ContextHit@K} =
\frac{1}{|\mathcal{Q}_{GT}|}
\sum_{q_t \in \mathcal{Q}_{GT}}
\mathbb{1}
\left[
R_{t,\mathrm{ctx}}^K \cap G_t \neq \varnothing
\right],
\]

and report context retention relative to dense retrieval:

\[
\mathrm{ContextRetention@K} =
\frac{\mathrm{ContextHit@K}}
     {\mathrm{DenseHit@K}}.
\]

If $\mathrm{DenseHit@K}=0$, we set $\mathrm{ContextRetention@K}=0$.

\subsection{Prequential Simulated Feedback}
\label{sec:05-experimental-setup:prequential-simulated-feedback}

LinUCB experiments use a no-leakage prequential simulated-feedback protocol. For each query, the current policy state is frozen before retrieval. The system ranks candidates, constructs the final context, and is evaluated against the ground-truth evidence. Only after this evaluation is the ground-truth label converted into simulated feedback and used to update the LinUCB state for later queries.

The feedback signal is controlled and ground-truth-derived. Oracle feedback is used only as an upper bound. Equal noisy and trust-weighted modes simulate imperfect user feedback with different reliability assumptions. Trust weighting changes how strongly a feedback event updates the route policy, but it does not give the current query access to its answer label before ranking.

This setup validates the route-learning mechanism under controlled feedback quality. It does not claim that real user feedback has already been collected, nor that delayed, biased, or adversarial human feedback would have the same effect without additional deployment safeguards.

Some experiments use multiple prequential epochs over the same query stream to simulate repeated interaction. These runs are useful for route-policy self-evolution analysis. They are not IID held-out generalization results.

\subsection{Implementation Notes}
\label{sec:05-experimental-setup:implementation-notes}

The main dense baseline uses \texttt{sentence-transformers/all-MiniLM-L6-v2} with exact cosine search on CPU. Embeddings and retrieval artifacts are cached to avoid repeating deterministic computation. Metrics are recomputed from saved rankings, not copied from prior summaries.

KMeans/MiniBatchKMeans uses a fixed number of arms across scales. This supports LinUCB comparability and reproducibility, but it is not claimed to be the best possible clustering design.
